% sec9_4.tex — Section 9.4: Augmented Dickey-Fuller Tests

The I(1) null process in the previous section was a random walk, which does not
allow for serial correlation in first differences. The tests of this section control for
serial correlation by adding lagged first differences to the autoregressive equation.
They are due to Dickey and Fuller (1979) and came to be known as the \textbf{Augmented
Dickey-Fuller} (\textbf{ADF}) \textbf{tests}.

\subsection*{The Augmented Autoregression}

The I(1) null in the simplest ADF tests is that $\{y_t\}$ is ARIMA($p$, 1, 0), that is, $\{\Delta y_t\}$
is a zero-mean stationary AR($p$) process:
\begin{equation}
\label{eq:9.4.1}
\Delta y_t = \zeta_1\,\Delta y_{t-1} + \zeta_2\,\Delta y_{t-2} + \cdots + \zeta_p\,\Delta y_{t-p} + \varepsilon_t,
\quad \E(\varepsilon_t^2) \equiv \sigma^2,
\end{equation}
where $\{\varepsilon_t\}$ is independent white noise and where the associated polynomial equation, $1 - \zeta_1 z - \zeta_2 z^2 - \cdots - \zeta_p z^p = 0$, has all its roots outside the unit circle. (Later
on we will allow $\{\Delta y_t\}$ to be stationary ARMA($p$, $q$).) The first difference process
$\{\Delta y_t\}$ is thus an I(0) process because it can be written as $\Delta y_t = \psi(L)\varepsilon_t$ where
\begin{equation}
\label{eq:9.4.2}
\psi(L) = (1 - \zeta_1 L - \zeta_2 L^2 - \cdots - \zeta_p L^p)^{-1}.
\end{equation}
The model used to derive the test statistics includes this process as a special case:
\begin{equation}
\label{eq:9.4.3}
y_t = \rho\, y_{t-1} + \zeta_1\,\Delta y_{t-1} + \zeta_2\,\Delta y_{t-2} + \cdots + \zeta_p\,\Delta y_{t-p} + \varepsilon_t.
\end{equation}
Indeed, \eqref{eq:9.4.3} reduces to~\eqref{eq:9.4.1} when $\rho = 1$.

Equation~\eqref{eq:9.4.3} can also be written as an AR($p+1$):
\begin{equation}
\label{eq:9.4.4}
y_t = \phi_1 y_{t-1} + \phi_2 y_{t-2} + \cdots + \phi_{p+1} y_{t-p-1} + \varepsilon_t,
\end{equation}
where the $\phi$'s are related to $(\rho, \zeta_1, \ldots, \zeta_p)$ by
\begin{align}
\rho &\equiv \phi_1 + \phi_2 + \cdots + \phi_{p+1}, \notag \\
\zeta_j &\equiv -(\phi_{j+1} + \phi_{j+2} + \cdots + \phi_{p+1})
\quad (j = 1, 2, \ldots, p).
\label{eq:9.4.5}
\end{align}
Conversely, any ($p+1$)-th order autoregression can be written in the form~\eqref{eq:9.4.3}.
This form~\eqref{eq:9.4.3} was originally proposed by Fuller (1976, p.~374). It is called the
\textbf{augmented autoregression} because it can be obtained by adding lagged changes
to the level AR(1) equation~\eqref{eq:9.3.1}.

It is easy to show (Review Question~2) that the polynomial equation associated
with~\eqref{eq:9.4.4} has a root inside the unit circle if $\rho > 1$. Thus, if we take the position
that the DGP is either I(1) or I(0), then $\rho$ cannot be greater than~1 and the I(0)
alternative is that $\rho < 1$.

\subsection*{Limiting Distribution of the OLS Estimator}

Having set up the model, we now turn to the task of deriving the limiting distribution of the OLS estimate of $\rho$ in the augmented autoregression~\eqref{eq:9.4.3} under the
I(1) null. Numerically the same estimate of $\rho$ can be obtained from the AR($p+1$)
equation~\eqref{eq:9.4.4} by calculating the sum of the OLS estimates of $(\phi_1, \phi_2, \ldots, \phi_{p+1})$
(recall that $\rho = \phi_1 + \cdots + \phi_{p+1}$). For the purpose of deriving the limiting distribution, however, the augmented autoregression is more convenient. Under the I(1)
null, all the $p+1$ regressors in~\eqref{eq:9.4.4} are driftless I(1), but it masks the fact, to
be exploited below and clear from~\eqref{eq:9.4.3}, that $p$ linear combinations of them are
zero-mean I(0).

To simplify the calculation, consider the special case of $p = 1$ (one lagged
change in $y_t$ or AR(2)), so that the augmented autoregression is
\begin{equation}
\label{eq:9.4.6}
y_t = \rho\, y_{t-1} + \zeta_1\,\Delta y_{t-1} + \varepsilon_t.
\end{equation}
We assume that the sample is $(y_{-1}, y_0, \ldots, y_T)$ so that the augmented autoregression can be estimated for $t = 1, 2, \ldots, T$ (or alternatively, the sample is
$(y_1, y_2, \ldots, y_T)$ but the summations in what follows are from $t = 3$ to $t = T$
rather than from $t = 1$ to $t = T$). Write this as
\[
y_t = \mathbf{x}_t' \bm{\beta} + \varepsilon_t
\]
with
\begin{equation}
\label{eq:9.4.7}
\underset{(2 \times 1)}{\mathbf{x}_t} = \begin{bmatrix} y_{t-1} \\ \Delta y_{t-1} \end{bmatrix}, \qquad
\underset{(2 \times 1)}{\bm{\beta}} = \begin{bmatrix} \rho \\ \zeta_1 \end{bmatrix}.
\end{equation}
If $\mathbf{b} \equiv (\hat{\rho}, \hat{\zeta}_1)'$ is the OLS estimator of the augmented autoregression coefficients
$\bm{\beta}\; (= (\rho, \zeta_1)')$, the sampling error is given by
\begin{equation}
\label{eq:9.4.8}
\mathbf{b} - \bm{\beta}
= \left(\sum_{t=1}^{T} \mathbf{x}_t \mathbf{x}_t'\right)^{-1} \sum_{t=1}^{T} \mathbf{x}_t \cdot \varepsilon_t.
\end{equation}
The individual terms in~\eqref{eq:9.4.8} are given by
\begin{align}
\mathbf{X}'\mathbf{X} &\equiv \sum_{t=1}^{T} \mathbf{x}_t \mathbf{x}_t'
= \begin{bmatrix}
\sum_{t=1}^{T}(y_{t-1})^2 & \sum_{t=1}^{T} y_{t-1}\,\Delta y_{t-1} \\[4pt]
\sum_{t=1}^{T} \Delta y_{t-1}\, y_{t-1} & \sum_{t=1}^{T} (\Delta y_{t-1})^2
\end{bmatrix}, \notag \\[6pt]
\sum_{t=1}^{T} \mathbf{x}_t \cdot \varepsilon_t
&= \begin{bmatrix}
\sum_{t=1}^{T} y_{t-1}\, \varepsilon_t \\[4pt]
\sum_{t=1}^{T} \Delta y_{t-1}\, \varepsilon_t
\end{bmatrix}.
\label{eq:9.4.9}
\end{align}

We can now apply the same trick we used for the time regression of Section~2.12. For some nonsingular matrix $\bm{\Upsilon}_T$ to be specified below, rewrite~\eqref{eq:9.4.8} as
\begin{equation}
\label{eq:9.4.10}
\bm{\Upsilon}_T \cdot (\mathbf{b} - \bm{\beta})
= \bm{\Upsilon}_T \begin{bmatrix} \hat{\rho} - \rho \\ \hat{\zeta}_1 - \zeta_1 \end{bmatrix}
= \mathbf{A}_T^{-1}\, \mathbf{c}_T,
\end{equation}
where
\begin{equation}
\label{eq:9.4.11}
\mathbf{A}_T \equiv \bm{\Upsilon}_T^{-1}\left[\sum_{t=1}^{T} \mathbf{x}_t \mathbf{x}_t'\right]\bm{\Upsilon}_T^{-1}, \quad
\mathbf{c}_T \equiv \bm{\Upsilon}_T^{-1} \sum_{t=1}^{T} \mathbf{x}_t \cdot \varepsilon_t.
\end{equation}

We seek a matrix $\bm{\Upsilon}_T$ such that $\bm{\Upsilon}_T \cdot (\mathbf{b} - \bm{\beta})$ has a nondegenerate limiting distribution under the I(1) null. This is accomplished by setting
\begin{equation}
\label{eq:9.4.12}
\bm{\Upsilon}_T = \begin{bmatrix} T & 0 \\ 0 & \sqrt{T} \end{bmatrix}.
\end{equation}
With this choice of the scaling matrix $\bm{\Upsilon}_T$, it follows from~\eqref{eq:9.4.9} and~\eqref{eq:9.4.11} that
\begin{align}
\mathbf{A}_T &= \begin{bmatrix}
\frac{1}{T^2}\sum_{t=1}^{T}(y_{t-1})^2 &
\frac{1}{\sqrt{T}}\frac{1}{T}\sum_{t=1}^{T} y_{t-1}\,\Delta y_{t-1} \\[6pt]
\frac{1}{\sqrt{T}}\frac{1}{T}\sum_{t=1}^{T}\Delta y_{t-1}\, y_{t-1} &
\frac{1}{T}\sum_{t=1}^{T}(\Delta y_{t-1})^2
\end{bmatrix}, \notag \\[6pt]
\mathbf{c}_T &= \begin{bmatrix}
\frac{1}{T}\sum_{t=1}^{T} y_{t-1}\,\varepsilon_t \\[4pt]
\frac{1}{\sqrt{T}}\sum_{t=1}^{T} \Delta y_{t-1}\,\varepsilon_t
\end{bmatrix}.
\label{eq:9.4.13}
\end{align}

Let us examine the elements of $\mathbf{A}_T$ and $\mathbf{c}_T$ and derive their limiting distributions
under the null of $\rho = 1$.

\textbf{The (1, 1) element of $\mathbf{A}_T$.} Since $\{y_t\}$ is I(1) without drift, it converges in distribution to $\lambda^2 \cdot \int W^2\,\mathrm{d}r$ by Proposition~9.2(a) where $\lambda^2 \equiv$ long-run variance of
$\{\Delta y_t\}$.

\textbf{The (2, 2) element of $\mathbf{A}_T$.} It converges in probability (and hence in distribution)
to $\gamma_0$ ($\equiv \Var(\Delta y_t)$) because $\{\Delta y_t\}$, being stationary AR(1), is ergodic stationary with mean zero.

\textbf{The off-diagonal elements of $\mathbf{A}_T$.} They are $1/\sqrt{T}$ times
\begin{equation}
\label{eq:9.4.14}
\frac{1}{T}\sum_{t=1}^{T} \Delta y_{t-1}\, y_{t-1},
\end{equation}
which is the average of the product of zero-mean I(0) and driftless I(1) variables. By Proposition~9.2(b), the similar product $(1/T)\sum_{t=1}^{T} \Delta y_t\, y_{t-1}$ converges in distribution to a random variable. It is easy to show (Review Question~3) that~\eqref{eq:9.4.14}, too, converges to a random variable. So $1/\sqrt{T}$ times~\eqref{eq:9.4.14},
which is the off-diagonal elements of $\mathbf{A}_T$, converges in probability (hence in
distribution) to~0.

Therefore, $\mathbf{A}_T$, the properly rescaled $\mathbf{X}'\mathbf{X}$ matrix, is asymptotically diagonal:
\[
\mathbf{A}_T \;\dto\;
\begin{bmatrix}
\lambda^2 \cdot \int_0^1 W(r)^2\,\mathrm{d}r & 0 \\
0 & \gamma_0
\end{bmatrix},
\]
so
\begin{equation}
\label{eq:9.4.15}
(\mathbf{A}_T)^{-1} \;\dto\;
\begin{bmatrix}
\big(\lambda^2 \cdot \int_0^1 W(r)^2\,\mathrm{d}r\big)^{-1} & 0 \\
0 & \gamma_0^{-1}
\end{bmatrix}.
\end{equation}

We now turn to the elements of $\mathbf{c}_T$.

\textbf{The first element of $\mathbf{c}_T$.} Using the Beveridge-Nelson decomposition, it can be
shown (Analytical Exercise~4) that
\begin{equation}
\label{eq:9.4.16}
\frac{1}{T}\sum_{t=1}^{T} y_{t-1}\,\varepsilon_t
\;\dto\; c_1 \equiv \tfrac{1}{2}\sigma^2 \cdot \psi(1) \cdot [W(1)^2 - 1],
\end{equation}
for a general zero-mean I(0) process
\begin{equation}
\label{eq:9.4.17}
\Delta y_t = \psi(L)\varepsilon_t.
\end{equation}
In the present case, $\Delta y_t = \zeta_1 \Delta y_{t-1} + \varepsilon_t$ under the null, so $\psi(L) = (1 - \zeta_1 L)^{-1}$
and
\begin{equation}
\label{eq:9.4.18}
\psi(1) = \frac{1}{1 - \zeta_1}.
\end{equation}

\textbf{Second element of $\mathbf{c}_T$.} It is easy to show (Review Question~4) that $\{\Delta y_{t-1}\varepsilon_t\}$ is a
stationary and ergodic martingale difference sequence, so
\begin{equation}
\label{eq:9.4.19}
\text{second element of } \mathbf{c}_T
= \frac{1}{\sqrt{T}}\sum_{t=1}^{T}\Delta y_{t-1}\,\varepsilon_t
\;\dto\; c_2 \sim N(0, \gamma_0 \cdot \sigma^2),
\end{equation}
where $\gamma_0 \equiv \Var(\Delta y_t)$ and $\sigma^2 \equiv \Var(\varepsilon_t)$.

Substituting \eqref{eq:9.4.15}--\eqref{eq:9.4.19} into~\eqref{eq:9.4.10} and noting that the true value of
$\rho$ is unity under the I(1) null, we obtain, for the OLS estimate of the augmented
autoregression,\footnote{Here, we are using Lemma~2.3(b) to claim that $\mathbf{A}_T^{-1}\mathbf{c}_T$ converges in distribution to $\mathbf{A}^{-1}\mathbf{c}$ where $\mathbf{A}$ and $\mathbf{c}$
are the limiting random variables of $\mathbf{A}_T$ and $\mathbf{c}_T$, respectively. Note that, unlike in situations encountered for the
stationary case, the convergence of $\mathbf{A}_T$ to $\mathbf{A}$ is in distribution, not in probability. We have proved in the text that
$\mathbf{A}_T$ converges in distribution to $\mathbf{A}$ and $\mathbf{c}_T$ to $\mathbf{c}$, but we have not shown that $(\mathbf{A}_T, \mathbf{c}_T)$ converges in distribution
jointly to $(\mathbf{A}, \mathbf{c})$, which is needed when using Lemma~2.3(b). However, by, e.g., Theorem~2.2 of Chan and Wei
(1988), the convergence is joint.}
\begin{align}
\bm{\Upsilon}_T \cdot (\mathbf{b} - \bm{\beta})
&= \begin{bmatrix} T \cdot (\hat{\rho} - 1) \\ \sqrt{T}\,(\hat{\zeta}_1 - \zeta_1) \end{bmatrix}
\;\dto\;
\begin{bmatrix}
\big(\lambda^2 \cdot \int_0^1 W(r)^2\,\mathrm{d}r\big)^{-1} & 0 \\
0 & \gamma_0^{-1}
\end{bmatrix}
\begin{bmatrix} c_1 \\ c_2 \end{bmatrix} \notag \\[6pt]
&= \begin{bmatrix}
(\lambda^2 \cdot \int_0^1 W(r)^2\,\mathrm{d}r)^{-1} c_1 \\
\gamma_0^{-1} c_2
\end{bmatrix}.
\label{eq:9.4.20result}
\end{align}
Thus,
\begin{align}
T \cdot (\hat{\rho} - 1) &\;\dto\;
\frac{\sigma^2 \psi(1)}{\lambda^2} \cdot \frac{\tfrac{1}{2}[W(1)^2 - 1]}{\int_0^1 W(r)^2\,\mathrm{d}r}
\quad\text{or}\quad
\frac{\lambda^2}{\sigma^2 \psi(1)} \cdot T \cdot (\hat{\rho} - 1) \;\dto\; DF_\rho,
\label{eq:9.4.20} \\
\sqrt{T}\,(\hat{\zeta}_1 - \zeta_1) &\;\dto\;
N\!\left(0,\;\frac{\sigma^2}{\gamma_0}\right),
\label{eq:9.4.21}
\end{align}
where $DF_\rho$ is none other than the DF $\rho$ distribution defined in~\eqref{eq:9.3.6}. Note that
the estimated coefficient of the zero-mean I(0) variable ($\Delta y_{t-1}$) has the usual $\sqrt{T}$
asymptotics, while that of the driftless I(1) variable ($y_{t-1}$) has a nonstandard limiting distribution.

\subsection*{Deriving Test Statistics}

There are no nuisance parameters entering the limiting distribution of the statistic
$\frac{\lambda^2}{\sigma^2 \psi(1)} \cdot T \cdot (\hat{\rho} - 1)$ in~\eqref{eq:9.4.20}, but the statistic itself involves nuisance parameters
through $\frac{\lambda^2}{\sigma^2 \psi(1)}$. By the formula~\eqref{eq:9.2.4} for the long-run variance $\lambda^2$ of $\{\Delta y_t\}$, this
ratio equals $\psi(1)$, which in turn equals $\frac{1}{1-\zeta_1}$ by~\eqref{eq:9.4.18} under the null. Now go
back to the augmented autoregression~\eqref{eq:9.4.6}. Since $\hat{\zeta}_1$, the OLS estimate of $\zeta_1$, is
consistent by~\eqref{eq:9.4.21}, a consistent estimator of $\frac{1}{1-\zeta_1}$ is $(1 - \hat{\zeta}_1)^{-1}$. It then follows
from~\eqref{eq:9.4.20} that
\begin{equation}
\label{eq:9.4.22}
\frac{T \cdot (\hat{\rho} - 1)}{1 - \hat{\zeta}_1} \;\dto\; DF_\rho.
\end{equation}
That is, the required correction on $T \cdot (\hat{\rho} - 1)$ is done through the estimated coefficient of lagged $\Delta y$ in the augmented autoregression. After the correction is made,
the limiting distribution of the statistic does not depend on nuisance parameters
controlling serial correlation of $\{\Delta y_t\}$ such as $\lambda^2$ and $\gamma_0$.

The OLS $t$-value for the hypothesis that $\rho = 1$, too, has a limiting distribution.
It can be written as
\begin{equation}
\label{eq:9.4.23}
t = \frac{\hat{\rho} - 1}{s \cdot \sqrt{(1, 1) \text{ element of } \big(\sum_{t=1}^{T} \mathbf{x}_t \mathbf{x}_t'\big)^{-1}}}
= \frac{T \cdot (\hat{\rho} - 1)}{s \cdot \sqrt{(1,1) \text{ element of } (\mathbf{A}_T)^{-1}}},
\end{equation}
where $s$ is the usual OLS standard error of regression. It follows easily from
\eqref{eq:9.4.15}, \eqref{eq:9.4.20}, the consistency of $s^2$ for $\sigma^2$, and Proposition~9.2(a) and~(b) for
$\xi_t = y_t$ that
\begin{equation}
\label{eq:9.4.24}
t \;\dto\; DF_t,
\end{equation}
where $DF_t$ is the DF $t$ distribution defined in~\eqref{eq:9.3.9}. Therefore, for the $t$-value,
there is no need to correct for the fact that lagged $\Delta y$ is included in the augmented
autoregression.

All these results for $p = 1$ can be generalized easily:

\begin{proposition}[Augmented Dickey-Fuller test of a unit root without intercept]
\label{prop:9.6}
Suppose that $\{y_t\}$ is ARIMA($p$, 1, 0) so that $\{\Delta y_t\}$ is a zero-mean stationary AR($p$) process following~\eqref{eq:9.4.1}. Consider estimating the augmented autoregression,~\eqref{eq:9.4.3}, and let $(\hat{\rho}, \hat{\zeta}_1, \hat{\zeta}_2, \ldots, \hat{\zeta}_p)$ be the OLS coefficient estimates. Also
let $t$ be the $t$-value for the hypothesis that $\rho = 1$. Then
\begin{align}
&\textit{ADF $\rho$ statistic:}\quad
\frac{T \cdot (\hat{\rho} - 1)}{1 - \hat{\zeta}_1 - \hat{\zeta}_2 - \cdots - \hat{\zeta}_p}
\;\dto\; DF_\rho,
\label{eq:9.4.25} \\
&\textit{ADF $t$ statistic:}\quad t \;\dto\; DF_t,
\label{eq:9.4.26}
\end{align}
where $DF_\rho$ and $DF_t$ are the two random variables defined in~\eqref{eq:9.3.6} and~\eqref{eq:9.3.9}.
\end{proposition}

\subsection*{Testing Hypotheses about $\zeta$}

The argument leading to this surprisingly simple result is lengthy, but it has a
by-product. The obvious generalization of~\eqref{eq:9.4.21} to the $p$-th order autoregression is
\begin{equation}
\label{eq:9.4.27}
\sqrt{T} \cdot
\begin{bmatrix}
\hat{\zeta}_1 - \zeta_1 \\
\hat{\zeta}_2 - \zeta_2 \\
\vdots \\
\hat{\zeta}_p - \zeta_p
\end{bmatrix}
\;\dto\;
N\!\left(
\begin{bmatrix} 0 \\ 0 \\ \vdots \\ 0 \end{bmatrix},\;
\sigma^2 \cdot
\begin{bmatrix}
\gamma_0 & \gamma_1 & \cdots & \gamma_{p-1} \\
\gamma_1 & \gamma_0 & \cdots & \gamma_{p-2} \\
\vdots & \vdots & \ddots & \vdots \\
\gamma_{p-1} & \gamma_{p-2} & \cdots & \gamma_0
\end{bmatrix}^{-1}
\right),
\end{equation}
where $\gamma_j$ is the $j$-th order autocovariance of $\{\Delta y_t\}$. By Proposition~6.7, this is the
same asymptotic distribution you would get if the augmented autoregression~\eqref{eq:9.4.3}
were estimated by OLS with the restriction $\rho = 1$, that is, if~\eqref{eq:9.4.1} is estimated
by OLS. It is easy to show (Review Question~6) that hypotheses involving only
the coefficients of the zero-mean I(0) regressors $(\zeta_1, \zeta_2, \ldots, \zeta_p)$ can be tested in the
usual way, with the usual $t$ and $F$ statistics asymptotically valid.

In deriving the limiting distributions of the regression coefficients, we have
exploited the fact that one of the regressors is zero-mean I(0) while the other is
driftless I(1). A systematic treatment of more general cases can be found in Sims,
Stock, and Watson (1990) and Watson (1994, Section~2).

\subsection*{What to Do When $p$ Is Unknown?}

Proposition~9.6 assumes that the order of autoregression, $p$, for $\Delta y_t$ is known.
What can we do when the order $p$ is unknown? To distinguish between the true
order $p$ and the number of lagged changes included in the augmented autoregression, we denote the latter by $\hat{p}$ in this section. We have considered a similar problem in Section~6.4. The difference is that the DGP here is a stationary process in
first differences, not in levels as in Section~6.4, and that the estimation equation is
an augmented autoregression with $\rho$ freely estimated. We display below, without
proof, three large-sample results about the choice of $\hat{p}$ under which the conclusion of Proposition~9.6 remains valid. These results are applicable to a class of
processes more general than is assumed in Proposition~9.6: $\{\Delta y_t\}$ is zero-mean
stationary and invertible ARMA($p$, $q$) of unknown order (albeit with an additional
assumption on $\varepsilon_t$ that its fourth moment is finite).\footnote{See Section~6.2 for the definition of invertible ARMA($p$, $q$) processes. If an ARMA($p$, $q$) process is invertible, then it can be written as an infinite-order autoregression.}

So, written as an autoregression, the order of autoregression for $\Delta y_t$ can be
infinite, as when $q > 0$, or finite, as when $q = 0$.

The first result is that the conclusion of Proposition~9.6 continues to hold when
$\hat{p}$, the number of lagged first differences in the augmented autoregression, is increased with the sample size at an appropriate rate.

\begin{enumerate}
\item (Said and Dickey, 1984, generalization of Proposition~9.6) \enspace Suppose that $\hat{p}$
satisfies
\begin{equation}
\label{eq:9.4.28}
\hat{p} \to \infty \quad \text{but} \quad \frac{\hat{p}}{T^{1/3}} \to 0 \quad \text{as } T \to \infty.
\end{equation}
(That is, $\hat{p}$ goes to infinity but at a rate slower than $T^{1/3}$.) Then the two
statistics, the ADF $\rho$ and ADF $t$, based on the augmented autoregression with
$\hat{p}$ lagged changes, have the same limiting distributions as those indicated in
Proposition~9.6.
\end{enumerate}

This result, however, does not give us a practical guide for the lag length selection because there are infinitely many rules for $\hat{p}$ satisfying~\eqref{eq:9.4.28}. A natural
question is whether any of the rules for selecting the lag length considered in
Section~6.4---the general-to-specific sequential $t$ rule or the information-criterion-based rules---can be used in the present context. To refresh your memory, the
information-criterion-based rule is to set $\hat{p}$ to the $j$ that minimizes
\begin{equation}
\label{eq:9.4.29}
\log\!\left(\frac{SSR_j}{T}\right) + (j + 1) \times \frac{C(T)}{T},
\end{equation}
where $SSR_j$ is the sum of squared residuals from the autoregression with $j$ lags:
\begin{equation}
\label{eq:9.4.30}
\Delta y_t = \rho\, y_{t-1} + \zeta_1\,\Delta y_{t-1} + \zeta_2\,\Delta y_{t-2} + \cdots + \zeta_j\,\Delta y_{t-j} + \varepsilon_t.
\end{equation}
The term $C(T)/T$ is multiplied by $j + 1$, because the equation has $j + 1$ coefficients including the lagged $y_t$ coefficient. For the Akaike information criterion
(AIC), $C(T) = 2$, while for the Bayesian information criterion (BIC, also called
Schwarz information criterion (SIC)), $C(T) = \log(T)$. In either rule, $\hat{p}$ is selected
from the set $j = 0, 1, 2, \ldots, p_{\max}$. In Section~6.4, this upper bound $p_{\max}$ was set
to some known integer greater than or equal to the true order of the finite-order
autoregressive process. The $\hat{p}$ selected by either of these rules is a function of data
(not just a function of $T$), and hence is a random variable.

To be sure, when $\{\Delta y_t\}$ is stationary and invertible ARMA($p$, $q$), the autoregressive order is infinite when $q > 0$ and the upper bound $p_{\max}$ obviously cannot
be set to the true order. But it can be made to increase with the sample size $T$. To
emphasize the dependence of $p_{\max}$ on $T$, write it as $p_{\max}(T)$. If the rules are thus
modified, does the Said-Dickey extension of Proposition~9.6 remain valid when $\hat{p}$
is chosen by either of the rules? The answer obtained by Ng and Perron (1995) is
yes. More specifically,

\begin{enumerate}
\setcounter{enumi}{1}
\item Suppose that $\hat{p}$ is selected by the general-to-specific $t$ rule with $p_{\max}(T)$ satisfying condition~\eqref{eq:9.4.28} (which in the Said-Dickey extension above was a
condition for $\hat{p}$, not for $p_{\max}(T)$) and $p_{\max}(T) > c \cdot T^g$ for some $c > 0$ and
$0 < g < 1/3$. Then the limiting distributions of the ADF $\rho$ and ADF $t$ statistics
are as indicated in Proposition~9.6.\footnote{This is Theorem~5.3 of Ng and Perron (1995). Their condition A1 is~\eqref{eq:9.4.28} here. Their condition A2$''$ is
implied by the second condition indicated here.}

\item Suppose that $\hat{p}$ is selected by AIC or by BIC with $p_{\max}(T)$ satisfying condition~\eqref{eq:9.4.28}. Then the same conclusion holds.\footnote{This is Theorem~4.3 of Ng and Perron (1995). The theorem does not indicate the upper bound $p_{\max}(T)$, but
(as pointed out by Pierre Perron in private communications with the author) Hannan and Deistler (1988, Section~7.4~(ii)) implies that~\eqref{eq:9.4.28} can be an upper bound.}
\end{enumerate}

Thus, no matter which rule---the sequential $t$, AIC, or BIC---for selecting $\hat{p}$ you
employ, the large-sample distributions of the ADF statistics are the same as in
Proposition~9.6.

\subsection*{A Suggestion for the Choice of $p_{\max}(T)$}

The finite-sample distribution, however, depends on which rule you use and also
on the choice of the upper bound $p_{\max}(T)$, and there are infinitely many valid
choices for the function $p_{\max}(T)$. For example, $p_{\max}(T) = [T^{1/4}]$ (the integer
part of $T^{1/4}$) satisfies the conditions in result~(2) for the sequential $t$-test. So does
$p_{\max}(T) = [100 T^{3/10}]$. It is important, therefore, that researchers use the same
$p_{\max}(T)$ when deciding the order of the augmented autoregression. In this context,
the Monte Carlo literature examining the small-sample properties of various unit-root tests is relevant. Simulations run by Schwert (1989) show that in small and
moderately large samples (from $T = 25$ to 1,000), including enough lags in the
autoregression is important to minimize size distortions.\footnote{Recall that a size distortion refers to the difference between the actual or exact size (the finite-sample probability of rejecting a true null) and the nominal size (the probability of rejecting a true null when the sample size
is infinite).} The choice of $\hat{p}$ (the
number of lags included in the autoregression) that was more or less successful in
controlling the actual size in Schwert's study is $\big[12 \cdot \big(\frac{T}{100}\big)^{1/4}\big]$. The upper bound
$p_{\max}(T)$ should therefore give lag selection rules a chance of selecting a $\hat{p}$ as large
as this. Therefore, in the examples and application of this chapter, we use the function
\begin{equation}
\label{eq:9.4.31}
p_{\max}(T) = \left[12 \cdot \left(\frac{T}{100}\right)^{1/4}\right]
\quad \text{(integer part of $12 \cdot (T/100)^{1/4}$)}
\end{equation}
in any of the rules for selecting the lag length. This function satisfies the conditions
of results~(2) and~(3) above.

The sample period when selecting $\hat{p}$ is $t = p_{\max}(T) + 2, \ldots, T$.
The first $t$ is $p_{\max}(T) + 2$ because $p_{\max}(T) + 1$ observations are needed to calculate $p_{\max}(T)$ lagged changes in the augmented autoregression. Since only $T -
p_{\max}(T) - 1$ observations are used to estimate the autoregression~\eqref{eq:9.4.30} for $j =
1, 2, \ldots, p_{\max}(T)$, the sample size in the objective function in~\eqref{eq:9.4.29} should be
$T - p_{\max}(T) - 1$ rather than $T$:
\begin{equation}
\label{eq:9.4.32}
\log\!\left(\frac{SSR_j}{T - p_{\max}(T) - 1}\right)
+ (j + 1) \times \frac{C(T - p_{\max}(T) - 1)}{T - p_{\max}(T) - 1}.
\end{equation}

\subsection*{Including the Intercept in the Regression}

As in the DF tests, we can modify the ADF tests so that they are invariant to an
addition of a constant to the series. We allow $\{y_t\}$ to differ by an unspecified
constant from the process that follows the augmented autoregression without intercept~\eqref{eq:9.4.3}. This amounts to replacing the $y_t$ in~\eqref{eq:9.4.3} by $y_t - \alpha$, which yields
\begin{equation}
\label{eq:9.4.33}
y_t = \alpha^* + \rho\, y_{t-1} + \zeta_1\,\Delta y_{t-1} + \zeta_2\,\Delta y_{t-2}
+ \cdots + \zeta_p\,\Delta y_{t-p} + \varepsilon_t,
\end{equation}
where $\alpha^* = (1 - \rho)\alpha$. As in the AR(1) model with intercept, the I(1) null is the
joint hypothesis that $\rho = 1$ \emph{and} $\alpha^* = 0$. Nevertheless, unit-root tests usually focus
on the single restriction $\rho = 1$.

It is left as an analytical exercise to prove

\begin{proposition}[Augmented Dickey-Fuller test of a unit root with intercept]
\label{prop:9.7}
Suppose that $\{y_t\}$ is an ARIMA($p$, 1, 0) process so that $\{\Delta y_t\}$ is a zero-mean
stationary AR($p$) process following~\eqref{eq:9.4.1}. Consider estimating the augmented
autoregression with intercept,~\eqref{eq:9.4.33}, and let $(\hat{\alpha}, \hat{\rho}^\mu, \hat{\zeta}_1, \hat{\zeta}_2, \ldots, \hat{\zeta}_p)$ be the OLS
coefficient estimates. Also let $t^\mu$ be the $t$-value for the hypothesis that $\rho = 1$. Then
\begin{align}
&\textit{ADF $\rho^\mu$ statistic:}\quad
\frac{T \cdot (\hat{\rho}^\mu - 1)}{1 - \hat{\zeta}_1 - \hat{\zeta}_2 - \cdots - \hat{\zeta}_p}
\;\dto\; DF_\rho^\mu,
\label{eq:9.4.34} \\
&\textit{ADF $t^\mu$ statistic:}\quad t^\mu \;\dto\; DF_t^\mu,
\label{eq:9.4.35}
\end{align}
where the random variables $DF_\rho^\mu$ and $DF_t^\mu$ are defined in~\eqref{eq:9.3.13} and~\eqref{eq:9.3.14}.
\end{proposition}

The basic idea of the proof is to use the Frisch-Waugh Theorem to write $\hat{\rho}^\mu$ and
the $t$-value in formulas involving the demeaned series created from $\{y_t\}$.

Before turning to an example, two comments about Proposition~9.7 are in order.

\begin{itemize}
\item (Numerical invariance) \enspace As was true in Proposition~9.4, since the regression
includes a constant, the test statistics are invariant to an addition of a constant to
the series.

\item (Said-Dickey extension) \enspace The Said-Dickey-Ng-Perron extension is also applicable here: if $\{\Delta y_t\}$ is stationary and invertible ARMA($p$, $q$) (so $\Delta y_t$ can be
written as a possibly infinite-order autoregression), then the ADF $\rho^\mu$ and $t^\mu$
statistics have the limiting distributions indicated in Proposition~9.7, provided
that $\hat{p}$ (the number of lagged changes to be included in the autoregression) is set
by the sequential $t$ rule or the AIC or BIC rules.\footnote{Ng and Perron (1995) proved this result about the selection of lag length only for the case without intercept.
That it can be extended to the case with intercept and also to the case with trend was confirmed in a private
communication with one of the authors of the paper.}
\end{itemize}

\begin{example}[ADF on T-bill rate]
\label{ex:9.2}
In the empirical exercise of Chapter~2, we remarked rather casually that the nominal interest rate $R_t$ might have
a trend for the sample period studied by Fama (1975). Here we test whether
it has a stochastic trend (i.e., whether it is driftless I(1)). The data set we
used in the empirical application of Chapter~2 was monthly observations on
the U.S.\ one-month Treasury bill rate. We take the sample period to be from
January 1953 to July 1971 ($T = 223$ observations), the sample period of
Fama's (1975) study. The interest rate series does not exhibit time trend. So
we include a constant but not time in the autoregression.

The maximum lag length $p_{\max}(T)$ is set equal to 14 according to the function~\eqref{eq:9.4.31}. In the process of choosing the actual lag length $\hat{p}$, we fix the
sample period from $t = p_{\max}(T) + 2 = 16$ (April 1954) to $t = 223$ (July
1971), with a sample size of $208$ ($= T - p_{\max}(T) - 1$). The sequential $t$ rule
picks a $\hat{p}$ of 14. The objective function in the information-criterion-based
rule when the sample size is fixed is~\eqref{eq:9.4.32}. The AIC picks $\hat{p} = 14$, while
the BIC picks $\hat{p} = 1$. Given $\hat{p} = 1$, we estimate the augmented autoregression on the maximum sample period, which is $t = \hat{p} + 2, \ldots, T$ (from March
1953 to July 1971) with 221 observations. The estimated regression is
\[
R_t = \underset{(0.057)}{0.088} + \underset{(0.0162)}{0.97705}\; R_{t-1}
- \underset{(0.067)}{0.20}\;\Delta R_{t-1},
\]
\[
R^2 = 0.944, \quad \text{sample size} = 221.
\]
From this estimated regression, we can calculate
\[
\text{ADF } \hat{\rho}^\mu = 221 \times \frac{0.97705 - 1}{1 + 0.20} = -4.22,
\quad t^\mu = \frac{0.97705 - 1}{0.0162} = -1.42.
\]
From Table~9.1, Panel~(b), the 5 percent critical value for the ADF $\rho^\mu$ statistic
is $-14.1$. The 5 percent critical value for the ADF $t^\mu$ statistic is $-2.86$ from
Table~9.2, Panel~(b). For either statistic, the I(1) null is easily accepted.
\end{example}

\subsection*{Incorporating the Time Trend}

Allowing for a linear trend, too, can be done as in the DF tests. Replacing the
$y_t$ in augmented autoregression~\eqref{eq:9.4.3} by $y_t - \alpha - \delta \cdot t$, we obtain the following
augmented autoregression with time trend:
\begin{equation}
\label{eq:9.4.36}
y_t = \alpha^* + \delta^* \cdot t + \rho\, y_{t-1} + \zeta_1\,\Delta y_{t-1}
+ \zeta_2\,\Delta y_{t-2} + \cdots + \zeta_p\,\Delta y_{t-p} + \varepsilon_t,
\end{equation}
where
\begin{equation}
\label{eq:9.4.37}
\alpha^* = (1-\rho)\alpha + (\rho - \zeta_1 - \zeta_2 - \cdots - \zeta_p)\delta,
\quad \delta^* = (1-\rho)\delta.
\end{equation}
Since $\delta^* = 0$ when $\rho = 1$, the I(1) null (that the process is I(1) with or without
drift) implies that $\rho = 1$ \emph{and} $\delta^* = 0$ in~\eqref{eq:9.4.36}, but, as usual, we focus on the
single restriction $\rho = 1$.

Combining the techniques used for proving Propositions~9.5 and~9.7, it is easy
to prove

\begin{proposition}[Augmented Dickey-Fuller Test of a unit root with linear time
trend]
\label{prop:9.8}
Suppose that $\{y_t\}$ is the sum of a linear time trend and an ARIMA($p$, 1, 0)
process so that $\{\Delta y_t\}$ is a stationary AR($p$) process whose mean may or may not
be zero. Consider estimating the augmented autoregression with trend,~\eqref{eq:9.4.36},
and let $(\hat{\alpha}, \hat{\delta}, \hat{\rho}^\tau, \hat{\zeta}_1, \hat{\zeta}_2, \ldots, \hat{\zeta}_p)$ be the OLS coefficient estimates. Also let $t^\tau$ be
the $t$-value for the hypothesis that $\rho = 1$. Then
\begin{align}
&\textit{ADF $\rho^\tau$ statistic:}\quad
\frac{T \cdot (\hat{\rho}^\tau - 1)}{1 - \hat{\zeta}_1 - \hat{\zeta}_2 - \cdots - \hat{\zeta}_p}
\;\dto\; DF_\rho^\tau,
\label{eq:9.4.38} \\
&\textit{ADF $t^\tau$ statistic:}\quad t^\tau \;\dto\; DF_t^\tau,
\label{eq:9.4.39}
\end{align}
where the random variables $DF_\rho^\tau$ and $DF_t^\tau$ are defined in~\eqref{eq:9.3.19} and~\eqref{eq:9.3.20}.
\end{proposition}

Before turning to an example, three comments are in order.

\begin{itemize}
\item (Numerical invariance to trend parameters) \enspace As in Proposition~9.5, since the
regression includes a constant and time, neither $\alpha$ nor $\delta$ affects the OLS estimates of $(\rho, \zeta_1, \ldots, \zeta_p)$ and their standard errors, so the finite-sample as well as
large-sample distribution of the ADF statistics will not be affected by $(\alpha, \delta)$.

\item (Said-Dickey extension) \enspace Again, the Said-Dickey-Ng-Perron extension is
applicable here: if $\{\Delta y_t\}$ is stationary and invertible ARMA($p$, $q$) with possibly
a nonzero mean, then the ADF $\rho^\tau$ and $t^\tau$ statistics have the limiting distributions
indicated in Proposition~9.8, provided that $\hat{p}$ is set by the sequential $t$-rule or the AIC or BIC rules.

\item (Should the regression include time?) \enspace The same comment we made about the
choice between the DF tests with or without trend is applicable to the ADF tests.
If you are confident that the null process has no trend, then the ADF $\rho^\mu$ and $t^\mu$
tests of Proposition~9.7 should be used, with the augmented autoregression not
including time, because the finite-sample power of the ADF tests is generally
greater without time in the augmented autoregression. On the other hand, if you
think that the process might have a trend, then you should use the ADF $\rho^\tau$ and
$t^\tau$ tests of Proposition~9.8 with time in the augmented autoregression.
\end{itemize}

\begin{example}[Does GDP have a unit root?]
\label{ex:9.3}
As was shown in Figure~9.1,
the log of U.S.\ GDP clearly has a linear time trend, so we include time in
the augmented autoregression. Using the logarithm of U.S.\ real GDP quarterly data from 1947:Q1 to 1998:Q1 (205 observations), we estimate the
augmented autoregression~\eqref{eq:9.4.36}. As in Example~9.2, we set $p_{\max}(T) =
\big[12 \cdot (205/100)^{1/4}\big]$ (which is 14). Also as in Example~9.2, the sample period
is fixed ($t = p_{\max}(T) + 2, \ldots, T$) in the process of selecting the number of
lagged changes. The sequential $t$ picks $\hat{p} = 12$, while both AIC and BIC
pick $\hat{p} = 1$. Given $\hat{p} = 1$, the augmented autoregression estimated on the
maximal sample of $t = \hat{p} + 2, \ldots, T$ (1947:Q3 to 1998:Q1) is
\[
y_t = \underset{(0.10)}{0.22} + \underset{(0.00011)}{0.00022}\;t
+ \underset{(0.0137)}{0.9707}\;y_{t-1}
+ \underset{(0.066)}{0.348}\;\Delta y_{t-1},
\]
\[
R^2 = 0.999, \quad \text{sample size} = 203.
\]
From this estimated regression, we can calculate
\begin{align*}
\text{ADF } \hat{\rho}^\tau &= 203 \times \frac{0.9707388 - 1}{1 - 0.348} = -9.11, \\
t^\tau &= \frac{0.9707388 - 1}{0.0137104} = -2.13.
\end{align*}
From Table~9.1, Panel~(c), the 5 percent critical value for the ADF $\rho^\tau$ statistic
is $-21.7$. The 5 percent critical value for the ADF $t^\tau$ statistic is $-3.41$ from
Table~9.2, Panel~(c). For either statistic, the I(1) null is easily accepted.

A little over 50 years of data may not be enough to discriminate between
the I(1) and I(0) alternatives. (As noted by Perron (1991), among others,
the span of the data rather than the sampling frequency, e.g., quarterly vs.\
annual, is more important for the power of unit-root tests, because $1 - \hat{p}$ is
much smaller on quarterly data than on annual data.) We now use annual
GDP data from 1869 to 1997 for the ADF tests.\footnote{Annual GDP data from 1929 are available from the Bureau of Economic Analysis web site. The Balke-Gordon GNP data (Balke and Gordon, 1989) were used to extrapolate GDP back to 1869.} So $T = 129$. As before,
we set $p_{\max}(T) = [12 \cdot (129/100)^{1/4}]$ (which is 12) and fix the sample period
in the process of choosing the lag length. The sequential $t$ picks 9 for $\hat{p}$,
while the value chosen by AIC and BIC is~1. When $\hat{p} = 1$, the regression
estimated on the maximum sample of $t = \hat{p} + 2, \ldots, T$ (from 1871 to 1997,
127 observations) is
\[
y_t = \underset{(0.18)}{0.67} + \underset{(0.0014)}{0.0048}\;t
+ \underset{(0.03989)}{0.85886}\;y_{t-1}
+ \underset{(0.086)}{0.32}\;\Delta y_{t-1},
\]
\[
R^2 = 0.999, \quad \text{sample size} = 127.
\]
The ADF $\rho^\tau$ and $t^\tau$ statistics are $-26.2$ and $-3.53$, respectively, which are
less (greater in absolute value) than their respective 5 percent critical values
of $-21.7$ and $-3.41$. With the annual data covering a much longer span of
time, we can reject the I(1) null that GDP has a stochastic trend.
\end{example}

The impression one gets from Examples~9.1--9.3 is the ease with which the I(1)
null is accepted, which leads to the suspicion that the DF and ADF tests may have
low power in finite samples. The power issue will be examined in the next section.

\subsection*{Summary of the DF and ADF Tests and Other Unit-Root Tests}

The various I(1) processes we have considered in this and the previous section
satisfy the model (a set of DGPs)
\begin{equation}
\label{eq:9.4.40}
y_t = d_t + z_t, \quad z_t = \rho z_{t-1} + u_t, \quad
\{u_t\} \sim \text{zero-mean I(0).}
\end{equation}
The restrictions on~\eqref{eq:9.4.40}, besides the restriction that $\rho = 1$, that characterize the
I(1) null for each of the cases considered are the following.

\begin{enumerate}
\item Proposition~9.3: $d_t = 0$, $\{u_t\}$ independent white noise,
\item Proposition~9.4: $d_t = \alpha$, $\{u_t\}$ independent white noise,
\item Proposition~9.5: $d_t = \alpha + \delta \cdot t$, $\{u_t\}$ independent white noise,
\item Said-Dickey extension of Proposition~9.6: $d_t = 0$, $\{u_t\}$ is zero-mean
ARMA($p$, $q$),
\item Said-Dickey extension of Proposition~9.7: $d_t = \alpha$, $\{u_t\}$ is zero-mean
ARMA($p$, $q$),
\item Said-Dickey extension of Proposition~9.8: $d_t = \alpha + \delta \cdot t$, $\{u_t\}$ is zero-mean
ARMA($p$, $q$).
\end{enumerate}

\bigskip
\noindent\textsc{questions for review} \hrulefill

\begin{enumerate}
\item (Non-uniqueness of the decomposition between zero-mean I(0) and driftless
I(1)) \enspace Assume $p = 2$ in~\eqref{eq:9.4.4}. Show that the third-order autoregression can
be written equivalently as
\[
y_t = a_1 \cdot y_{t-1} + a_2 \cdot (y_{t-1} - y_{t-2}) + a_3 \cdot (y_{t-1} - y_{t-3}) + \varepsilon_t.
\]
How are $(a_1, a_2, a_3)$ related to $(\phi_1, \phi_2, \phi_3)$? Which regressor is zero-mean I(0)
and which one is driftless I(1)? Verify that the OLS estimate of the $y_{t-1}$ coefficient from this equation is numerically the same as the OLS estimate of $\rho$ from~\eqref{eq:9.4.3} with $p = 2$.

\item Let $\phi(z) = 1 - \phi_1 z - \cdots - \phi_{p+1} z^{p+1}$ be the polynomial associated with the
AR($p+1$) process~\eqref{eq:9.4.4}. Show that $\phi(z) = 0$ has a real root inside the unit
circle if $\rho > 1$. \textbf{Hint:} $\phi(1) = 1 - (\phi_1 + \cdots + \phi_{p+1}) = 1 - \rho < 0$ if $\rho > 1$.
$\phi(0) = 1 > 0$.

\item Prove that~\eqref{eq:9.4.14} $\dto$ $\frac{\lambda^2}{2} W(1)^2 + \frac{\gamma_0}{2}$. \textbf{Hint:} The asymptotic distribution
of $\frac{1}{T}\sum_{t=1}^{T}\Delta y_{t-1}\, y_{t-1}$ is the same as that of $\frac{1}{T}\sum_{t=1}^{T}\Delta y_t\, y_t$. Also, $\Delta y_t\, y_t =
\Delta y_t\, y_{t-1} + (\Delta y_t)^2$. Use Proposition~9.2(b).

\item Let $\{\Delta y_t\}$ be a zero-mean I(0) process satisfying \eqref{eq:9.2.1}--\eqref{eq:9.2.3b}. (The zero-mean stationary AR($p$) process $\{\Delta y_t\}$ considered in Proposition~9.6 is a
special case of this.)

\begin{enumerate}[label=(\alph*)]
\item Show that $\{\Delta y_{t-1}\varepsilon_t\}$ is a stationary and ergodic martingale difference
sequence. \textbf{Hint:} Since $\{\varepsilon_t\}$ and $\{\Delta y_t\}$ are jointly ergodic stationary, $\{\Delta y_{t-1} \cdot
\varepsilon_t\}$ is ergodic stationary. To show that $\{\Delta y_{t-1}\varepsilon_t\}$ is a martingale difference
sequence, use the Law of Iterated Expectations. $\{\varepsilon_{t-1}, \varepsilon_{t-2}, \ldots\}$ has more
information than $\{\Delta y_{t-2}, \Delta y_{t-3}, \varepsilon_{t-2}, \ldots\}$.

\item Prove~\eqref{eq:9.4.19}. \textbf{Hint:} You need to show that $\E[(\Delta y_{t-1}\,\varepsilon_t)^2] = \gamma_0 \cdot \sigma^2$.
\end{enumerate}

\item Verify~\eqref{eq:9.4.24}. \textbf{Hint:} From~\eqref{eq:9.4.20} and~\eqref{eq:9.2.4},
\[
T \cdot (\hat{\rho} - 1) \;\dto\;
\frac{1}{\psi(1)} \cdot \frac{\tfrac{1}{2}(W(1)^2 - 1)}{\int_0^1 W^2\,\mathrm{d}r}.
\]
Also, by~\eqref{eq:9.4.15} and~\eqref{eq:9.2.4},
\begin{equation}
\label{eq:9.4.41}
s \cdot \sqrt{(1,1) \text{ element of } (\mathbf{A}_T)^{-1}}
\;\dto\; 1 \Big/ \sqrt{\psi(1)^2 \int_0^1 W^2\,\mathrm{d}r}.
\end{equation}
Since $\psi(1) > 0$ by the stationarity condition, $\sqrt{\psi(1)^2} = \psi(1)$.

\item ($t$-test on $\zeta$) \enspace It is claimed in the text (right below~\eqref{eq:9.4.27}) that the usual $t$-value for each $\zeta$ is asymptotically standard normal. Verify this for the case of
$p = 1$. \textbf{Hint:}~\eqref{eq:9.4.27} for $p = 1$ is
\[
\sqrt{T} \cdot (\hat{\zeta}_1 - \zeta_1) \;\dto\; N(0, \sigma^2/\gamma_0).
\]
You need to show that $\sqrt{T}$ times the standard error of $\hat{\zeta}_1$ converges in probability to the square root of $\sigma^2/\gamma_0$. $\sqrt{T}$ times the standard error of $\hat{\zeta}_1$ is $s \cdot
\sqrt{(2, 2) \text{ element of } (\mathbf{A}_T)^{-1}}$.
\end{enumerate}
